% 深圳大学实验报告 - Project 1
% 上传Overleaf后请选择XeLaTeX编译。
\documentclass[a4paper,11pt]{ctexart}
\setCJKmainfont{SimSun}[BoldFont=SimSun]
\setCJKsansfont{SimSun}[BoldFont=SimSun]
\setCJKmonofont{SimSun}
\usepackage[left=1.8cm,right=1.8cm,top=1.8cm,bottom=1.8cm]{geometry}
\usepackage{booktabs,array,xcolor,hyperref,fancyhdr,tikz,listings}
\usetikzlibrary{positioning}
\hypersetup{colorlinks=true,urlcolor=blue}
\definecolor{SZUBlue}{HTML}{005BAC}
\definecolor{CodeBg}{HTML}{F7F9FC}
\pagestyle{fancy}
\fancyhf{}
\lhead{深圳大学实验报告 | Python程序设计}
\rfoot{第 \thepage 页}
\setlength{\parindent}{2em}
\setlength{\parskip}{0.25em}
\lstset{basicstyle=\ttfamily\small,backgroundcolor=\color{CodeBg},frame=single,
  breaklines=true,columns=fullflexible}

\begin{document}

\begin{titlepage}
\centering
\vspace*{1.2cm}
{\Huge\bfseries 深\quad 圳\quad 大\quad 学\quad 实\quad 验\quad 报\quad 告\par}
\vspace{1.2cm}
{\Large\color{SZUBlue}\bfseries Project 1: NumPy图像编码器\par}
\vspace{0.4cm}
{\small OpenCLIP ViT-B-32 Image Encoder Forward in NumPy\par}
\vspace{1.4cm}
\renewcommand{\arraystretch}{1.8}
\begin{tabular}{>{\bfseries}p{3cm}p{9cm}}
课程名称 & Python程序设计\\\hline
项目名称 & NumPy图像编码器\\\hline
学院 & 人工智能学院\\\hline
专业 & 人工智能\\\hline
指导教师 & 樊超\\\hline
报告人 & 郑心而\\\hline
学号 & 2025300053\\\hline
实验时间 & 2026.11.17 - 2026.11.18\\\hline
提交时间 & 2026.11.18\\\hline
\end{tabular}
\vfill
\textbf{代码仓库：}\url{https://github.com/Miya555orz/my-numpy-homework}
\vspace{1cm}
\par 教务处制
\end{titlepage}

\section{实验目标}
只使用NumPy完成OpenCLIP ViT-B-32图像编码器前向计算，将输入
$[B,3,224,224]$ 转换为每张图的512维单位向量，并在给定容差内匹配独立FP32参考。

\section{方法与张量维度}
\begin{center}
\begin{tikzpicture}[node distance=0.7cm, every node/.style={draw=SZUBlue,rounded corners,fill=blue!5,minimum height=1cm,font=\small}]
\node (a) {图像 $[B,3,224,224]$};
\node[right=of a] (b) {Patch $[B,49,3072]$};
\node[right=of b] (c) {Token $[B,50,768]$};
\node[below=of c] (d) {12层Block};
\node[left=of d] (e) {CLS $[B,768]$};
\node[left=of e] (f) {输出 $[B,512]$};
\draw[->] (a)--(b);\draw[->] (b)--(c);\draw[->] (c)--(d);
\draw[->] (d)--(e);\draw[->] (e)--(f);
\end{tikzpicture}
\end{center}

图像块构造过程中的张量维度变化如下：
\[
[B,3,224,224]\rightarrow[B,3,7,32,7,32]
\rightarrow[B,7,7,3,32,32]\rightarrow[B,49,3072].
\]
轴置换使图像块按行优先顺序排列，并保持块内通道、高度和宽度的展开顺序。

\section{关键实现}
多头注意力将768维特征划分为12个64维注意力头，分别计算缩放点积注意力，
随后合并各头输出。每个Transformer模块包含两次残差连接，前馈网络的中间维度为3072。
\begin{lstlisting}[language=Python]
k_transpose = k.transpose(0, 1, 3, 2)
scores = (q @ k_transpose) / np.sqrt(np.float32(head_dim))
context = attention_weights(scores) @ v
merged = context.transpose(0, 2, 1, 3).reshape(B, N, 768)
\end{lstlisting}

\subsection{维度错误及修正}
对四维键张量使用二维示例中的\texttt{k.T}会反转全部轴，
导致批次维度与注意力头维度错位。因此实现采用\texttt{transpose(0,1,3,2)}，
仅交换最后两个轴，使注意力分数的维度为$[B,H,N,N]$。
注意力与前馈网络的线性权重采用$[out,in]$布局，计算时需要转置；
最终投影矩阵\texttt{visual.proj}采用$[in,out]$布局，计算时保持原方向。

\newpage
\section{实验结果}
\begin{center}
\begin{tabular}{lccc}
\toprule
阶段 & 输出shape & 最大绝对误差 & 平均绝对误差\\
\midrule
preprocess & (10,3,224,224) & 0.000e+00 & 0.000e+00\\
patches & (2,49,3072) & 0.000e+00 & 0.000e+00\\
tokens & (2,50,768) & 4.005e-05 & 1.289e-06\\
block0 & (2,50,768) & 7.063e-06 & 3.250e-07\\
last\_tokens & (2,50,768) & 4.768e-05 & 2.696e-06\\
output & (10,512) & 5.439e-07 & 8.311e-08\\
batch & (3,512) & 4.619e-07 & 7.632e-08\\
\bottomrule
\end{tabular}
\end{center}
各阶段数值检查均通过。最终输出的最大绝对误差为$5.439\times10^{-7}$，
小于该阶段的绝对容差$2\times10^{-5}$。乱序及重复样本的批次检查通过，
表明输出能够对应输入样本的实际顺序。

\section{讨论与总结}
张量维度管理是实现正确性的关键。图像块构造改变数据排列方式，
多头注意力对特征轴进行拆分与重组，残差连接则保持模块输入与输出维度一致。
逐阶段比较可定位误差首次出现的位置。实现中的主要性能开销来自给定的精确GELU函数，
其逐元素标量计算限制了运行速度；该限制不影响数值检查结果。

\section{问题修正与工作说明}
复核练习合并文件时，发现E5的函数名为\texttt{unit\_vector\_exercise}，
与独立题目要求的\texttt{unit\_vector}接口不一致。现已统一函数名及调用处，
并重新运行练习代码。M1的独立版本原先未拒绝零查询向量，
现已在分母计算前加入长度检查并抛出\texttt{ValueError}。
本人负责练习实现、编码器代码核对和结果验证。
报告整理使用课程提供的starter、检查器及OpenAI Codex辅助；
数值结果以本地检查器运行输出为依据。

\vfill
\noindent\textbf{复现命令：}\texttt{python check.py --stage all}。

\end{document}
